% 教材：Shankar, Principles of Quantum Mechanics, Second Edition
% 已覆盖：第 3 章第 3.1--3.5 节，印刷页 pp. 107--113，PDF pp. 120--126
% 人工核验：已完成分块讲解、问答和自编收口检查

\section{第 3.1 节：经典物理中的粒子与波}
\label{section:QM-C03-S01}

\studymeta
  {QM-C03-S01}
  {2026-09-04}
  {Shankar, \textit{Principles of Quantum Mechanics}, Second Edition}
  {第 3 章 \S 3.1；印刷页 pp.~107--108；PDF pp.~120--121}
  {讲解、自编检查与订正已完成}

\textbf{本节主线。}

经典物理把物质与辐射分别放入两套清楚分开的图景：粒子由有限个位置和动量变量描述，并沿确定轨迹运动；波则是定义在整个空间上的场，由波动方程控制。平面波把波长、周期、传播方向和强度压缩进振幅、频率与波矢。第三章将以这组经典区分为出发点，考察实验为何迫使我们修改它。

\notetopic{QM-C03-S01-T01}{经典粒子、经典波与初始数据}

\keyidea{粒子的状态由有限个数给出；波在固定时刻的状态是一整个空间函数。二者都需要与时间二阶运动方程相匹配的两组初始数据，但数据的类型不同。}

在经典粒子图景中，单个三维粒子在时刻 $t_0$ 的状态可由
\begin{equation}
  \bm r(t_0),\qquad \dot{\bm r}(t_0)
  \label{eq:QM-C03-S01-particle-initial-data}
\end{equation}
描述；在 Hamiltonian 形式中也可改用 $\bm r(t_0),\bm p(t_0)$。给定运动方程和初值后，原则上可确定一条轨迹 $\bm r(t)$。粒子的能量与动量被看作由局域对象携带；对有限个粒子，状态只需有限个自由度。

经典波则是分布在空间中的扰动。例如一维弦的位移由
\[
  \psi(x,t)
\]
描述：固定 $t$ 后，$x$ 是标记各空间位置的参数，而 $\psi(x,t)$ 给出每一点的场值。连续介质因此具有连续无穷多个自由度，不能用一个粒子坐标 $x(t)$ 代替整幅波形。

典型三维波动方程为
\begin{equation}
  \boxed{
  \nabla^2\psi(\bm r,t)
  =\frac{1}{c^2}
  \frac{\partial^2\psi(\bm r,t)}{\partial t^2}.}
  \label{eq:QM-C03-S01-wave-equation}
\end{equation}
由于它对时间是二阶方程，确定后续演化通常需要两个空间函数：
\begin{equation}
  \psi(\bm r,t_0),
  \qquad
  \left.\frac{\partial\psi(\bm r,t)}{\partial t}\right|_{t=t_0}.
  \label{eq:QM-C03-S01-wave-initial-data}
\end{equation}
它们分别是初始波形和每个空间点的初始变化率，对应粒子问题中的初始位置与初始速度。

只给初始波形一般不能确定传播方向。以一维波为例，
\begin{equation}
  \psi_+(x,t)=f(x-ct),
  \qquad
  \psi_-(x,t)=f(x+ct)
  \label{eq:QM-C03-S01-opposite-traveling-waves}
\end{equation}
在 $t=0$ 时都有 $\psi_\pm(x,0)=f(x)$，但
\[
  \dot\psi_+(x,0)=-cf'(x),
  \qquad
  \dot\psi_-(x,0)=cf'(x).
\]
前者向右、后者向左传播。相同初始形状配上不同初始速度，可以产生不同的后续运动。

\notetopic{QM-C03-S01-T02}{一维平面波、相位与传播速度}

一维简谐平面波可写成
\begin{align}
  \psi(x,t)
  &=A\exp\left[\ii\left(
    \frac{2\pi}{\lambda}x-
    \frac{2\pi}{T}t
  \right)\right] \notag\\
  &=\boxed{A\ee^{\ii(kx-\omega t)}},
  \label{eq:QM-C03-S01-plane-wave}
\end{align}
其中
\begin{equation}
  k=\frac{2\pi}{\lambda},
  \qquad
  \omega=\frac{2\pi}{T}.
  \label{eq:QM-C03-S01-k-omega-definitions}
\end{equation}
$k$ 是 wavenumber（波数或空间角频率），$\omega$ 是 angular frequency（角频率）。空间前进一个波长 $\lambda$ 或时间经过一个周期 $T$，相位都恰好改变 $2\pi$；因此 $k$ 与 $\omega$ 分别比 $1/\lambda$ 与 $1/T$ 多一个 $2\pi$。

波的相位为
\[
  \phi(x,t)=kx-\omega t.
\]
追踪某个固定相位点，令 $kx-\omega t=\phi_0$，则
\begin{equation}
  x(t)=\frac{\omega}{k}t+\frac{\phi_0}{k},
  \qquad
  v_{\mathrm{ph}}=\frac{\dd x}{\dd t}=\frac{\omega}{k}
  =\frac{\lambda}{T}.
  \label{eq:QM-C03-S01-phase-velocity}
\end{equation}
在 $k,\omega>0$ 的约定下，$\ee^{\ii(kx-\omega t)}$ 向 $+x$ 传播，而 $\ee^{\ii(kx+\omega t)}$ 向 $-x$ 传播。

把式~\eqref{eq:QM-C03-S01-plane-wave} 代入一维波动方程，可以直接得到频率与波数的关系：
\begin{align}
  \frac{\partial^2\psi}{\partial x^2}
  &=-k^2\psi,
  &
  \frac{\partial^2\psi}{\partial t^2}
  &=-\omega^2\psi.
  \label{eq:QM-C03-S01-plane-wave-second-derivatives}
\end{align}
因此
\begin{equation}
  -k^2\psi
  =-\frac{\omega^2}{c^2}\psi
  \quad\Longrightarrow\quad
  \boxed{\omega^2=c^2k^2}.
  \label{eq:QM-C03-S01-dispersion-relation}
\end{equation}
这个 $\omega$ 与 $k$ 的函数关系称为 dispersion relation（色散关系）。对正频率、向右传播的分支，$\omega=ck$，所以相速度就是 $c$。

\dialoguesupplement{与标准形式 $A\ee^{\ii(kx-\omega t)}$ 比较时，时间项的系数是 $-\omega$。例如 $3\ee^{\ii(4x-12t)}$ 对应 $k=4$、$\omega=12$，而不是 $\omega=-12$；周期和速率均取正值，传播方向由固定相位随时间的移动方向判断。}

\notetopic{QM-C03-S01-T03}{复数表示、强度与三维波矢}

真实的经典场通常取实值。若
\[
  A=\abs{A}\ee^{\ii\alpha},
\]
则复平面波的实部为
\begin{equation}
  \operatorname{Re}\psi(x,t)
  =\abs{A}\cos(kx-\omega t+\alpha),
  \label{eq:QM-C03-S01-real-wave}
\end{equation}
其中 $\alpha$ 是初相位。复指数是计算工具：微分只会产生 $\ii k$ 或 $-\ii\omega$ 因子。由于实系数波动方程是线性的，若复函数 $\psi=\psi_R+\ii\psi_I$ 是解，则实部 $\psi_R$ 与虚部 $\psi_I$ 分别也是解；计算结束后可取所需的实部。

教材采用
\begin{equation}
  I=\abs{\psi}^2
  \label{eq:QM-C03-S01-intensity}
\end{equation}
表示强度。对理想平面波，
\begin{equation}
  \abs{\psi(x,t)}^2
  =\abs{A}^2
   \abs{\ee^{\ii(kx-\omega t)}}^2
  =\boxed{\abs{A}^2},
  \label{eq:QM-C03-S01-plane-wave-intensity}
\end{equation}
因为复相位只在单位圆上转动而不改变模长。实际经典波的时间平均能流与振幅平方成正比；具体比例系数取决于所描述的物理场和归一化约定。

三维平面波写成
\begin{equation}
  \boxed{
  \psi(\bm r,t)
  =A\ee^{\ii(\bm k\cdot\bm r-\omega t)}},
  \qquad
  \abs{\bm k}=\frac{2\pi}{\lambda}.
  \label{eq:QM-C03-S01-three-dimensional-wave}
\end{equation}
$\bm k$ 是 wave vector（波矢）：它的模长给出空间相位变化率，方向给出传播方向。在固定时刻，等相位面满足
\begin{equation}
  \bm k\cdot\bm r=\text{常数}.
  \label{eq:QM-C03-S01-wavefront}
\end{equation}
因此 $\bm k$ 是等相位面的法向量，而单位法向量是
\[
  \widehat{\bm k}=\frac{\bm k}{\abs{\bm k}}.
\]
等相位面与 $\bm k$ 垂直，而不是平行。

\begin{figure}[htbp]
  \centering
  \begin{tikzpicture}[x=1.05cm,y=1.05cm,>=Latex]
    \foreach \x in {-2,-1,0,1,2}
      \draw[noteBlue,thick] (\x,-1.35)--(\x,1.35);
    \draw[->,very thick,accentRed] (-2.55,0)--(2.65,0)
      node[right] {$\bm k$};
    \node[noteBlue,align=center] at (0,1.75)
      {等相位面在二维截面中的投影};
    \draw[<->,verifyOrange] (0,-1.65)--(1,-1.65)
      node[midway,below] {$\lambda$};
  \end{tikzpicture}
  \caption{相邻波面相隔一个波长；波沿波矢所指的法向传播。}
  \label{fig:QM-C03-S01-wavefronts}
\end{figure}

三维相速度的大小为
\begin{equation}
  v_{\mathrm{ph}}=\frac{\omega}{\abs{\bm k}}.
  \label{eq:QM-C03-S01-three-dimensional-phase-velocity}
\end{equation}
本节只讨论单一平面波的相位传播；不同波数分量叠加后的波包及其传播将在后文需要时再区分处理。

\textbf{一分钟回顾。}

经典粒子由有限个位置与动量变量描述，经典波则由整个空间上的场函数描述。波动方程对时间为二阶，因此初始波形与初始变化率缺一不可。一维平面波 $A\ee^{\ii(kx-\omega t)}$ 的固定相位以 $\omega/k$ 向右传播，并满足 $\omega^2=c^2k^2$。复指数简化微分运算，实际场可取其实部；理想平面波强度为 $\abs{A}^2$。三维波矢 $\bm k$ 指向传播方向，并垂直于等相位面。

\section{第 3.2 节：经典波与经典粒子的双缝实验}
\label{section:QM-C03-S02}

\studymeta
  {QM-C03-S02}
  {2026-09-05}
  {Shankar, \textit{Principles of Quantum Mechanics}, Second Edition}
  {第 3 章 \S 3.2；印刷页 pp.~108--110；PDF pp.~121--123}
  {讲解、自编检查与订正已完成}

\textbf{本节主线。}

同一套双缝装置把经典波与经典粒子的组合规则清楚地区分开来：经典波先叠加复振幅，再取模平方得到强度，因此会出现干涉交叉项；彼此独立的经典粒子沿互斥路径到达探测器，因而直接相加概率或计数率。教材将在下一节用这两个经典基准检验光。

\notetopic{QM-C03-S02-T01}{经典波：振幅叠加、干涉项与条纹位置}

设两缝间距为 $d$，缝屏到探测屏的距离为 $D$，屏上观察点为 $X$，其纵坐标为 $x$；从两缝 $S_1,S_2$ 到 $X$ 的路程分别为 $r_1,r_2$。符号约定如图~\ref{fig:QM-C03-S02-double-slit-geometry} 所示。

\begin{figure}[htbp]
  \centering
  \begin{tikzpicture}[x=1.08cm,y=1.08cm,>=Latex]
    \foreach \y in {-1.5,-0.9,-0.3,0.3,0.9,1.5}
      \draw[->,noteBlue] (-2.0,\y)--(-0.25,\y);

    \draw[line width=1.8pt] (0,-2.0)--(0,-0.95);
    \draw[line width=1.8pt] (0,-0.55)--(0,0.55);
    \draw[line width=1.8pt] (0,0.95)--(0,2.0);
    \fill[accentRed] (0,0.75) circle (1.5pt)
      node[left=3pt] {$S_1$};
    \fill[accentRed] (0,-0.75) circle (1.5pt)
      node[left=3pt] {$S_2$};

    \draw[line width=1.3pt] (5,-2.0)--(5,2.0)
      node[above] {探测屏};
    \fill[verifyOrange] (5,1.35) circle (1.7pt)
      node[right=3pt] {$X$};
    \coordinate (O) at (5,0);
    \draw[densely dashed,noteBlue] (0,0.75)--(5,1.35)
      node[pos=0.57,above] {$r_1$};
    \draw[densely dashed,noteBlue] (0,-0.75)--(5,1.35)
      node[pos=0.57,below] {$r_2$};

    \draw[<->,verifyOrange] (0.38,-0.75)--(0.38,0.75)
      node[midway,right] {$d$};
    \draw[<->] (0,-2.25)--(5,-2.25)
      node[midway,below] {$D$};
    \draw[<->] (5.35,0)--(5.35,1.35)
      node[midway,right] {$x$};
    \draw[densely dotted] (0,0)--(O);
  \end{tikzpicture}
  \caption{双缝几何示意图。图为笔记自绘；远屏近轴近似要求 $D\gg d$ 且 $\abs{x}\ll D$。}
  \label{fig:QM-C03-S02-double-slit-geometry}
\end{figure}

只开 $S_1$ 时，屏上的复振幅和强度分别记作
\[
  \phi_1(x,t),
  \qquad
  I_1(x)=\abs{\phi_1(x,t)}^2;
\]
只开 $S_2$ 时同理得到 $\phi_2$ 与 $I_2$。两缝同时打开时，线性波动方程要求先叠加振幅：
\begin{equation}
  \phi_{1+2}=\phi_1+\phi_2.
  \label{eq:QM-C03-S02-amplitude-addition}
\end{equation}
所以强度为
\begin{align}
  I_{1+2}
  &=\abs{\phi_1+\phi_2}^2 \notag\\
  &=(\phi_1+\phi_2)^*(\phi_1+\phi_2) \notag\\
  &=\abs{\phi_1}^2+\abs{\phi_2}^2
    +\phi_1^*\phi_2+\phi_2^*\phi_1 \notag\\
  &=\boxed{I_1+I_2+2\operatorname{Re}(\phi_1^*\phi_2)}.
  \label{eq:QM-C03-S02-wave-interference}
\end{align}
最后两项是 interference term（干涉项）。这里必须保留复共轭；一般不能把干涉项写成 $\phi_1\phi_2$。

若两条波等振幅且同频，取
\begin{equation}
  \phi_j=A\ee^{\ii(kr_j-\omega t)},
  \qquad
  I_0=\abs{A}^2,
  \qquad
  \Delta r=r_2-r_1,
  \label{eq:QM-C03-S02-equal-amplitude-waves}
\end{equation}
则
\begin{align}
  I_{1+2}
  &=2I_0+2I_0\cos(k\Delta r) \notag\\
  &=\boxed{4I_0\cos^2\left(\frac{k\Delta r}{2}\right)}.
  \label{eq:QM-C03-S02-interference-intensity}
\end{align}
因此
\begin{align}
  \Delta r=n\lambda
  &\quad\Longrightarrow\quad
  I_{\max}=4I_0,
  \label{eq:QM-C03-S02-bright-condition}\\
  \Delta r=\left(n+\frac12\right)\lambda
  &\quad\Longrightarrow\quad
  I_{\min}=0,
  \label{eq:QM-C03-S02-dark-condition}
\end{align}
其中 $n\in\mathbb Z$。明纹对应同相叠加，暗纹对应反相抵消。

在 $D\gg d$ 且 $\abs{x}\ll D$ 的远屏近轴条件下，
\begin{equation}
  \Delta r\simeq d\sin\theta\simeq\frac{dx}{D}.
  \label{eq:QM-C03-S02-path-difference-approximation}
\end{equation}
代入明暗条件可得
\begin{align}
  x_n^{\mathrm{bright}}
  &\simeq \frac{n\lambda D}{d},
  &
  x_n^{\mathrm{dark}}
  &\simeq \frac{(n+\tfrac12)\lambda D}{d},
  \label{eq:QM-C03-S02-fringe-positions}
\end{align}
相邻明纹或暗纹的间距均为
\begin{equation}
  \boxed{\Delta x\simeq\frac{\lambda D}{d}}.
  \label{eq:QM-C03-S02-fringe-spacing}
\end{equation}
因此增大波长或屏距会拉宽条纹，增大缝距会压缩条纹。

\notetopic{QM-C03-S02-T02}{经典粒子：互斥路径的概率直接相加}

把入射波换成具有固定能量、方向有一定分布的经典粒子束。令 $P_1(x)$ 与 $P_2(x)$ 表示在相同入射条件下，粒子分别经 $S_1$、$S_2$ 落入 $x$ 附近同一探测区间的概率贡献；它们不是已经分别归一化为 $1$ 的条件分布。若探测的是单位时间内的粒子数，也可用相应计数率 $I_1(x),I_2(x)$ 表示相同的分布规律。

对每个经典粒子而言，“经 $S_1$ 到达 $x$”与“经 $S_2$ 到达 $x$”是 mutually exclusive（互斥）的两条路径。在两缝之间没有额外相互作用、粒子束条件保持不变的经典假设下，概率加法法则给出
\begin{equation}
  \boxed{P_{1+2}(x)=P_1(x)+P_2(x)},
  \qquad
  \boxed{I_{1+2}(x)=I_1(x)+I_2(x)}.
  \label{eq:QM-C03-S02-classical-particle-addition}
\end{equation}
这里没有类似式~\eqref{eq:QM-C03-S02-wave-interference} 的交叉项。$P_1P_2$ 描述两个事件同时发生时才可能出现的乘法结构，并不适用于“经 $S_1$ 或经 $S_2$”这一互斥选择。

即使两缝发出的粒子可能彼此碰撞，也可把入射强度降到一次只有一个粒子通过装置，从而排除粒子之间的相互影响。单次事件仍是某一个探测器的一次局域计数；重复足够多次后形成直方图，其平滑包络趋近 $I_1$、$I_2$ 或 $I_{1+2}$。在经典粒子模型中，由于 $P_2(x)\geq 0$，打开第二条路径不可能使某处的总到达概率低于只开第一条路径时的概率。

\notetopic{QM-C03-S02-T03}{实验判据与相干性边界}

两套经典预测的差别不在探测器是否逐点计数，而在组合规则：
\begin{equation}
  \boxed{
  \begin{aligned}
    \text{经典波：}&\quad
    \text{先加振幅，再取模平方};\\
    \text{经典粒子：}&\quad
    \text{对互斥路径直接加概率}.
  \end{aligned}}
  \label{eq:QM-C03-S02-two-composition-rules}
\end{equation}
经典波可以满足
\[
  I_{1+2}(x)\neq I_1(x)+I_2(x).
\]
特别是在干涉极小处，打开第二条缝反而可能降低该处的能流。经典粒子则必须满足式~\eqref{eq:QM-C03-S02-classical-particle-addition}，因此这种局部降低不可能出现。这是比“分布看起来有起伏”更直接的实验判据。

\dialoguesupplement{看不到稳定条纹并不足以单独证明对象服从经典粒子加法。若两束经典波的相对相位 $\delta$ 随时间或样本随机变化，则
\[
  I_{1+2}=I_1+I_2+2\sqrt{I_1I_2}\cos\delta.
\]
在相位均匀漂移且观测进行时间平均或系综平均时，$\langle\cos\delta\rangle=0$，于是 $\langle I_{1+2}\rangle=I_1+I_2$。瞬时干涉仍可能存在，只是缺乏稳定的相干关系，平均后不再形成可见条纹。}

\newpage
\textbf{一分钟回顾。}

双缝实验中，经典波先叠加振幅，再计算强度：
\[
  \phi_{1+2}=\phi_1+\phi_2,
  \qquad
  I_{1+2}=I_1+I_2+\phi_1^*\phi_2+\phi_2^*\phi_1.
\]
等振幅时强度为 $4I_0\cos^2(k\Delta r/2)$，远屏近轴下条纹间距为 $\lambda D/d$。经典粒子的两条路径互斥，故 $P_{1+2}=P_1+P_2$，不存在干涉项。打开第二缝能否降低某处的探测率，是区分两种经典组合规则的关键；但“没有稳定条纹”还可能来自经典波失去相干性，不能单独作为粒子性的充分证据。

\section{第 3.3 节：光的双缝实验}
\label{section:QM-C03-S03}

\studymeta
  {QM-C03-S03}
  {2026-09-06}
  {Shankar, \textit{Principles of Quantum Mechanics}, Second Edition}
  {第 3 章 \S 3.3；印刷页 pp.~110--112；PDF pp.~123--125}
  {讲解、自编检查与订正已完成}

\textbf{本节主线。}

光在双缝实验中同时表现出两类不能由同一套经典图景解释的现象：探测器每次记录一个局域、不可继续分割的能量转移事件，而大量这类事件累积出的空间分布却服从振幅叠加产生的干涉规律。量子理论因此不把光简单归入“经典粒子”或“经典波”，而是以概率振幅描述单次事件发生在各处的可能性。

\notetopic{QM-C03-S03-T01}{光的离散探测与光子能量、动量}

在强光下，探测器接收的能量看起来近似连续，经典电磁波的强度描述十分有效。把光强逐步降低并使用足够灵敏的探测器后，能量交换却表现为一次次局域计数，而不是任意小份额的连续沉积。每次计数所对应的电磁辐射量子称为 photon（光子）。

对角频率为 $\omega$、普通频率为 $\nu$ 的单色光，单个光子的能量由实验规律给出：
\begin{equation}
  \boxed{E=\hbar\omega=h\nu},
  \qquad
  \hbar=\frac{h}{2\pi},
  \qquad
  \omega=2\pi\nu.
  \label{eq:QM-C03-S03-photon-energy}
\end{equation}
光子的动量与波矢满足
\begin{equation}
  \boxed{\bm p=\hbar\bm k},
  \qquad
  p=\hbar k=\frac{h}{\lambda}.
  \label{eq:QM-C03-S03-photon-momentum}
\end{equation}
这里的 $E=\hbar\omega$ 与 $\bm p=\hbar\bm k$ 是由实验建立的量子关系，并不能只从经典波动方程推出。经典理论给出的 $\omega$、$\bm k$ 在量子描述中进一步决定单个光子的能量和动量。

真空中的光满足色散关系 $\omega=ck$，所以
\begin{equation}
  E=\hbar\omega
  =\hbar ck
  =pc.
  \label{eq:QM-C03-S03-photon-dispersion}
\end{equation}
与相对论能量--动量关系
\begin{equation}
  E^2=p^2c^2+m_0^2c^4
  \label{eq:QM-C03-S03-relativistic-energy-momentum}
\end{equation}
比较可知光子的静质量 $m_0=0$。这并不意味着光子没有能量或动量；它们分别由式~\eqref{eq:QM-C03-S03-photon-energy} 和式~\eqref{eq:QM-C03-S03-photon-momentum} 给出。

在理想单色光和固定探测效率的条件下，若平均每秒到达的光子数为 $\dot N$，平均功率为
\begin{equation}
  \boxed{\mathcal P=\dot N\hbar\omega}.
  \label{eq:QM-C03-S03-power-count-rate}
\end{equation}
因此在频率不变时降低光强，不会把每个光子的能量按比例削小，而会降低平均光子到达率。例如光强降为原来的 $1/5$ 时，单个光子的 $\hbar\omega$ 不变，单位时间内的平均计数降为原来的 $1/5$。若频率改变，则单个光子的能量也随 $\omega$ 改变。

\notetopic{QM-C03-S03-T02}{单光子双缝实验为何排除经典轨迹模型}

把入射光调弱到装置中平均一次只有一个光子，探测屏仍会记录一个个局域事件。短时间内这些点看似随机；经过大量重复后，计数直方图逐渐显现与经典波强度相同的干涉条纹。由于同一时刻没有第二个光子通过装置，条纹不能解释为不同光子之间碰撞或彼此扰动的结果。

经典粒子模型要得到
\begin{equation}
  N_{1+2}(x)=N_1(x)+N_2(x),
  \label{eq:QM-C03-S03-classical-count-addition}
\end{equation}
实际上需要同时采用两项假设：每个粒子经过 $S_1$ 或 $S_2$ 中一条确定路径；并且打开另一条缝不会改变原来那条路径对 $x$ 处计数的贡献。在这些假设下，两类事件互斥，故概率或平均计数率只能相加。

单光子实验却一般满足
\begin{equation}
  \boxed{\rho_{1+2}(x)\neq\rho_1(x)+\rho_2(x)},
  \label{eq:QM-C03-S03-nonclassical-probability-addition}
\end{equation}
其中 $\rho$ 表示单位长度内的探测概率密度。最直接的反例是理想暗纹位置 $x_0$。若两条路径在该处分别贡献
\begin{equation}
  \psi_1(x_0)=A,
  \qquad
  \psi_2(x_0)=-A,
  \label{eq:QM-C03-S03-dark-fringe-amplitudes}
\end{equation}
则只开任一条缝时的概率密度均为 $\abs{A}^2$，而两缝同时打开时
\begin{equation}
  \rho_{1+2}(x_0)
  =\abs{A-A}^2=0.
  \label{eq:QM-C03-S03-dark-fringe-probability}
\end{equation}
打开第二条缝反而把该处的探测概率从正值降为零，这与式~\eqref{eq:QM-C03-S03-classical-count-addition} 冲突。

这里排除的是“光子总有一条确定经典路径，而且另一条缝不影响该路径贡献”的独立轨迹模型。仅凭本节实验，不应把结论夸大为排除了所有可能含有轨迹概念的量子诠释。教材在此要确立的操作性事实是：实验结果不能由经典互斥概率直接相加得到，而必须先组合复概率振幅。

\notetopic{QM-C03-S03-T03}{概率振幅、Born 规则与统计条纹}

用态矢量 $\ket{\psi}$ 描述入射光的量子态。在连续位置基 $\{\ket{x}\}$ 中，
\begin{equation}
  \boxed{\psi(x)=\braket{x}{\psi}}
  \label{eq:QM-C03-S03-position-amplitude}
\end{equation}
是态矢量在 $\ket{x}$ 方向上的坐标，称为在位置 $x$ 探测到光子的 probability amplitude（概率振幅）。它本身一般是复数，不是概率，也不是经典电磁场能量密度。

Born 规则把振幅与可观测概率联系起来：
\begin{align}
  \Pr(x\leq X<x+\dd x)
  &=\abs{\psi(x)}^2\dd x,
  \label{eq:QM-C03-S03-born-infinitesimal}\\
  \Pr(a\leq X\leq b)
  &=\int_a^b\abs{\psi(x)}^2\dd x.
  \label{eq:QM-C03-S03-born-interval}
\end{align}
若光子必在整条探测屏上被发现，则态须满足归一化条件
\begin{equation}
  \int_{-\infty}^{\infty}\abs{\psi(x)}^2\dd x=1.
  \label{eq:QM-C03-S03-normalization}
\end{equation}
因此 $\abs{\psi(x)}^2$ 的量纲是长度的倒数，$\psi(x)$ 的量纲是 $\mathrm{length}^{-1/2}$。单点在连续空间中的概率通常为零；有物理意义的是有限区间概率或概率密度。

两缝同时打开时，量子理论使用与经典波相同形式的振幅叠加规则：
\begin{equation}
  \boxed{\psi_{1+2}(x)=\psi_1(x)+\psi_2(x)}.
  \label{eq:QM-C03-S03-amplitude-superposition}
\end{equation}
但这里的 $\psi_j$ 是到达 $x$ 的概率振幅。应用 Born 规则后，
\begin{align}
  \rho_{1+2}(x)
  &=\abs{\psi_1(x)+\psi_2(x)}^2 \notag\\
  &=\abs{\psi_1(x)}^2+\abs{\psi_2(x)}^2
    +\psi_1^*(x)\psi_2(x)+\psi_2^*(x)\psi_1(x) \notag\\
  &=\boxed{\rho_1(x)+\rho_2(x)
    +2\operatorname{Re}\!\left[\psi_1^*(x)\psi_2(x)\right]}.
  \label{eq:QM-C03-S03-born-interference}
\end{align}
交叉项可以为正也可以为负，所以打开第二条缝既可能增加、也可能降低某处的探测概率。

一次实验只产生一个完整、局域的探测事件；式~\eqref{eq:QM-C03-S03-born-interference} 预测的是这种事件在大量相同制备下的空间分布。若重复发送 $N$ 个同样制备的光子，则区间 $[a,b]$ 内计数的期望值为
\begin{equation}
  \boxed{
  \mathbb E\!\left[N_{[a,b]}\right]
  =N\int_a^b\abs{\psi(x)}^2\dd x.}
  \label{eq:QM-C03-S03-expected-counts}
\end{equation}
随着样本数增加，归一化直方图才逐渐逼近 $\abs{\psi(x)}^2$。因此“一个光子的一次探测”与“许多次探测形成干涉条纹”处在不同统计层次，并不矛盾。

\dialoguesupplement{为什么不直接把光视为经典波？经典波理论能准确给出许多实验中的平均干涉包络，但它把场能量视为可连续分割，不能单独解释探测器为何以局域的、能量为 $\hbar\omega$ 的完整事件响应。量子理论保留振幅的叠加结构，同时用 Born 规则把振幅模平方解释为单次事件的概率密度。}

\begin{figure}[htbp]
  \centering
  \begin{tikzpicture}[
    >=Latex,
    node distance=0.72cm and 0.65cm,
    box/.style={draw=noteBlue,rounded corners=2pt,align=center,
      minimum height=0.82cm,text width=2.55cm,fill=noteBlue!5},
    event/.style={draw=accentRed,rounded corners=2pt,align=center,
      minimum height=0.82cm,text width=2.35cm,fill=accentRed!5}
  ]
    \node[box] (state) {同一制备的\\单光子态 $\ket{\psi}$};
    \node[box,right=of state] (paths) {两条路径的振幅贡献\\$\psi_1(x),\ \psi_2(x)$};
    \node[box,right=of paths] (sum) {总振幅\\$\psi_1(x)+\psi_2(x)$};
    \node[event,right=of sum] (click) {一次完整的\\局域探测事件};
    \draw[->,thick] (state)--(paths);
    \draw[->,thick] (paths)--(sum);
    \draw[->,thick] (sum)--node[above] {$\abs{\cdot}^2$} (click);

    \node[box,below=1.0cm of sum] (repeat) {重复 $N$ 次\\独立制备与探测};
    \node[event,right=of repeat] (hist) {计数直方图\\$\propto\abs{\psi(x)}^2$};
    \draw[->,thick] (click)--(repeat);
    \draw[->,thick] (repeat)--(hist);
  \end{tikzpicture}
  \caption{单光子振幅叠加与统计条纹的层次。图为笔记自绘；两条振幅贡献不是把一个光子的能量分成两半。}
  \label{fig:QM-C03-S03-single-photon-statistics}
\end{figure}

\textbf{术语与适用边界。}
所谓“一个光子与自身干涉”只是对“同一单光子态的不同路径振幅在探测前相干叠加”的简称，不表示光子彼此碰撞，也不表示一次探测只吸收半个光子。更现代的表述把光子视为量子化电磁场的一次激发；本节的单光子概率波图景是进入量子公设前的物理动机，并不是完整的量子电动力学。严格的光子位置表象还存在超出本节范围的技术问题。

\textbf{一分钟回顾。}

光以局域、离散事件与物质交换能量和动量，每个光子满足 $E=\hbar\omega$、$\bm p=\hbar\bm k$，且真空中 $E=pc$。在固定频率下减弱光强，主要降低的是平均光子到达率，而非单个光子的能量。即使一次只有一个光子通过双缝，累计结果仍出现干涉，说明经典互斥路径概率不能直接相加。量子理论先叠加概率振幅 $\psi_1+\psi_2$，再由 Born 规则取模平方得到探测概率密度；单次结果是一个完整局域事件，大量重复才显现 $\abs{\psi(x)}^2$ 所规定的统计条纹。

\section{第 3.4 节：物质波（de Broglie 波）}
\label{section:QM-C03-S04}

\studymeta
  {QM-C03-S04}
  {2026-09-07}
  {Shankar, \textit{Principles of Quantum Mechanics}, Second Edition}
  {第 3 章 \S 3.4；印刷页 p.~112；PDF p.~125}
  {讲解、教材估算例、自编检查与订正已完成}

\textbf{本节主线。}

光原本被视为波，却以光子的形式携带离散的能量和动量。de Broglie 反向追问：传统上被视为粒子的电子等物质，是否也具有波动性？他的假设把粒子动量与波矢联系起来，并预言物质粒子同样可以发生干涉和衍射。宏观物体并非没有物质波，而是其波长和条纹间距通常小到无法由现实仪器分辨。

\notetopic{QM-C03-S04-T01}{de Broglie 假设与物质波的概率解释}

de Broglie hypothesis（德布罗意假设）认为，动量为 $\bm p$ 的粒子对应波矢 $\bm k$，满足
\begin{equation}
  \boxed{\bm p=\hbar\bm k}.
  \label{eq:QM-C03-S04-de-broglie-vector}
\end{equation}
因此 $\bm k$ 与 $\bm p$ 同向，而且
\begin{equation}
  k=\frac{p}{\hbar}.
  \label{eq:QM-C03-S04-wave-number}
\end{equation}
结合 $k=2\pi/\lambda$ 和 $h=2\pi\hbar$，得到 de Broglie wavelength（德布罗意波长）
\begin{align}
  \frac{2\pi}{\lambda}
  &=\frac{p}{\hbar},
  \notag\\
  \boxed{\lambda=\frac{2\pi\hbar}{p}=\frac{h}{p}}.
  \label{eq:QM-C03-S04-de-broglie-wavelength}
\end{align}
真正普遍的关系是 $\lambda=h/p$。只有在非相对论条件下，才可再用 $p=mv$ 写成
\begin{equation}
  \boxed{\lambda=\frac{h}{mv}}.
  \label{eq:QM-C03-S04-nonrelativistic-wavelength}
\end{equation}
速度接近光速时必须使用相对论动量，不能继续把 $p$ 简写为 $mv$。

式~\eqref{eq:QM-C03-S04-de-broglie-wavelength} 给出两个直接判断：
\begin{itemize}[leftmargin=*]
  \item 对同一种粒子，动量越大，波长越短；
  \item 两个粒子若动量大小相同，则 de Broglie 波长相同；若非相对论速度相同，则质量越大的粒子波长越短。
\end{itemize}

Davisson--Germer experiment（戴维孙--革末实验）观察到了电子衍射，从实验上验证了电子的波数与动量服从 $k=p/\hbar$。物质波并不表示电子沿波浪形轨迹前进，也不应被解释为连续摊开的经典物质密度。教材沿用上一节的概率解释，以
\begin{equation}
  \psi(x,t)
  \label{eq:QM-C03-S04-matter-wave-amplitude}
\end{equation}
表示在时刻 $t$ 于位置 $x$ 探测到粒子的 probability amplitude（概率振幅），而
\begin{equation}
  \boxed{
  \Pr(x\leq X<x+\dd x)
  =\abs{\psi(x,t)}^2\dd x}
  \label{eq:QM-C03-S04-matter-wave-born-rule}
\end{equation}
才是相应的位置概率。

在双缝实验中，两条路径贡献的概率振幅先相加：
\begin{equation}
  \psi_{1+2}(x)=\psi_1(x)+\psi_2(x),
  \label{eq:QM-C03-S04-matter-wave-superposition}
\end{equation}
于是
\begin{equation}
  \abs{\psi_{1+2}(x)}^2
  =\abs{\psi_1(x)}^2+\abs{\psi_2(x)}^2
   +2\operatorname{Re}\!\left[\psi_1^*(x)\psi_2(x)\right].
  \label{eq:QM-C03-S04-matter-wave-interference}
\end{equation}
单个粒子仍只在屏上产生一次局域探测；大量同样制备的粒子才累积出由式~\eqref{eq:QM-C03-S04-matter-wave-interference} 描述的干涉分布。

\dialoguesupplement{本节所排除的是朴素经典模型：粒子经过一条确定且不受另一条缝影响的路径，两条互斥路径的概率直接相加。仅凭这里的论证，不应把结论扩大为排除一切含“轨迹”概念的量子诠释。}

\notetopic{QM-C03-S04-T02}{宏观经典极限：极短波长与不可分辨条纹}

教材考虑质量为 $1\,\mathrm{g}$、速度为 $1\,\mathrm{cm/s}$ 的弹丸。在非相对论条件下，
\begin{equation}
  p=mv
  =1\,\mathrm{g\,cm/s}.
  \label{eq:QM-C03-S04-pellet-momentum}
\end{equation}
使用
\[
  h\simeq
  6.626\times10^{-27}\,
  \mathrm{g\,cm^2/s},
\]
得到
\begin{align}
  \lambda
  &=\frac{h}{p}
    =\frac{6.626\times10^{-27}\,
      \mathrm{g\,cm^2/s}}
      {1\,\mathrm{g\,cm/s}}
  \notag\\
  &=6.626\times10^{-27}\,\mathrm{cm}
    \sim\boxed{10^{-26}\,\mathrm{cm}}.
  \label{eq:QM-C03-S04-pellet-wavelength}
\end{align}
教材指出，这比质子半径还小约 $10^{13}$ 倍。如此短的波长并不等于“没有物质波”，而意味着相应的干涉结构极其细密。

对远屏近轴双缝装置，条纹间距近似为
\begin{equation}
  \boxed{\Delta x\simeq\frac{\lambda D}{d}},
  \label{eq:QM-C03-S04-fringe-spacing}
\end{equation}
其中 $D$ 是双缝到屏幕的距离，$d$ 是缝间距。$\lambda$ 越短，$\Delta x$ 也越小。

实际探测器具有有限的空间分辨率。若一个像素覆盖宽度 $\delta x$，它记录的是区间概率
\begin{equation}
  P_{\mathrm{pixel}}
  =\int_{x_0}^{x_0+\delta x}
   \abs{\psi_{1+2}(x)}^2\dd x,
  \label{eq:QM-C03-S04-pixel-probability}
\end{equation}
而不是数学上无限精确的点值。当
\begin{equation}
  \delta x\gg\Delta x
  \label{eq:QM-C03-S04-unresolved-condition}
\end{equation}
时，一个像素内部包含许多明暗条纹。若 $\abs{\psi_1}$、$\abs{\psi_2}$ 的包络在该像素内变化缓慢，快速振荡的交叉项会近似平均为零：
\begin{equation}
  \int_{x_0}^{x_0+\delta x}
  2\operatorname{Re}\!\left[\psi_1^*(x)\psi_2(x)\right]\dd x
  \simeq0.
  \label{eq:QM-C03-S04-cross-term-average}
\end{equation}
因此仪器看到的粗粒化分布近似满足
\begin{equation}
  \boxed{
  \overline{\rho}_{1+2}(x)
  \simeq\rho_1(x)+\rho_2(x),}
  \label{eq:QM-C03-S04-classical-coarse-graining}
\end{equation}
与经典粒子预测相同。这里消失的是仪器可分辨的干涉条纹，而不是物质波关系本身。

\begin{figure}[htbp]
  \centering
  \begin{tikzpicture}[>=Latex,node distance=0.8cm and 0.7cm,
    box/.style={draw=noteBlue,rounded corners=2pt,align=center,
      minimum height=0.9cm,text width=2.55cm,fill=noteBlue!5},
    result/.style={draw=accentRed,rounded corners=2pt,align=center,
      minimum height=0.9cm,text width=2.7cm,fill=accentRed!5}]
    \node[box] (p) {宏观动量 $p$ 很大};
    \node[box,right=of p] (lambda) {$\lambda=h/p$\\极小};
    \node[box,right=of lambda] (fringe) {$\Delta x\simeq\lambda D/d$\\极小};
    \node[result,right=of fringe] (average) {有限分辨率平均\\干涉交叉项};
    \draw[->,thick] (p)--(lambda);
    \draw[->,thick] (lambda)--(fringe);
    \draw[->,thick] (fringe)--(average);
  \end{tikzpicture}
  \caption{本节给出的宏观经典极限。图为笔记自绘；经典分布来自不可分辨条纹的粗粒化平均，而非先验删除物质波。}
  \label{fig:QM-C03-S04-macroscopic-limit}
\end{figure}

\textbf{适用边界。}
上述“条纹过密而被平均”的解释就是教材本节采用的论证。现实宏观物体还会与环境强烈耦合，decoherence（退相干）也会抑制可观测干涉；这属于后续量子理论内容，不是本节推导的前提。

\textbf{一分钟回顾。}

de Broglie 假设把物质粒子的动量与波矢联系为 $\bm p=\hbar\bm k$，因而 $\lambda=h/p$；非相对论时才可进一步写成 $h/(mv)$。物质波 $\psi(x,t)$ 是概率振幅，而不是经典物质密度或波浪形轨迹。单个粒子产生局域探测，大量重复实验通过振幅交叉项形成干涉分布。宏观物体的动量很大，故 $\lambda$ 与双缝条纹间距都极小；当仪器分辨率远大于条纹间距时，交叉项被粗粒化平均，最终呈现近似经典的概率相加规律。

\section{第 3.5 节：结论}
\label{section:QM-C03-S05}

\studymeta
  {QM-C03-S05}
  {2026-09-08}
  {Shankar, \textit{Principles of Quantum Mechanics}, Second Edition}
  {第 3 章 \S 3.5；印刷页 pp.~112--113；PDF pp.~125--126}
  {讲解、自编检查与订正已完成}

\textbf{本节主线。}

第三章并未给出完整的量子理论，而是用双缝实验说明经典粒子与经典波两套图景都不足以单独描述电子、光子等量子对象。量子对象在探测时呈现局域、完整的事件，却在重复实验的统计分布中呈现概率振幅的叠加与干涉。由此需要用波函数或抽象态矢量描述状态，并把状态随时间的变化作为量子动力学的核心问题；这些经验事实将在下一章被组织成量子力学公设。

\notetopic{QM-C03-S05-T01}{波粒二象性的操作性内容}

量子对象的 particlelike（粒子性）首先体现在单次探测：电子或光子在屏幕上的某个确定位置触发一次局域事件，并在该事件中传递完整的能量、动量与电荷。一次结果可以写成
\begin{equation}
  X=x_{\mathrm{det}},
  \label{eq:QM-C03-S05-local-detection}
\end{equation}
但“探测位置确定”并不推出探测前存在一条确定的经典轨迹。

其 wavelike（波动性）则体现在概率振幅的组合规律。若到达同一位置 $x$ 的两种不可区分选择分别贡献 $\psi_1(x,t)$ 与 $\psi_2(x,t)$，则
\begin{align}
  \psi(x,t)
  &=\psi_1(x,t)+\psi_2(x,t),
  \label{eq:QM-C03-S05-amplitude-sum}\\
  \rho(x,t)
  &=\abs{\psi(x,t)}^2 \notag\\
  &=\abs{\psi_1(x,t)}^2+\abs{\psi_2(x,t)}^2
    +2\operatorname{Re}\!\left[\psi_1^*(x,t)\psi_2(x,t)\right].
  \label{eq:QM-C03-S05-interference-density}
\end{align}
交叉项使累计分布不能由两条经典互斥路径的概率直接相加得到。对单个粒子，位置测量的 Born 规则为
\begin{equation}
  \boxed{
  \Pr(a\leq X\leq b;t)
  =\int_a^b\abs{\psi(x,t)}^2\dd x.}
  \label{eq:QM-C03-S05-position-born-rule}
\end{equation}
因此 $\psi(x,t)$ 与每个粒子的量子态相关，而不是只有许多粒子组成的束流才拥有的经典波。一次实验给出一个位置；大量相同制备的实验才使归一化计数分布趋近 $\abs{\psi(x,t)}^2$。

wave--particle duality（波粒二象性）不应理解为量子对象在“经典小球”和“经典连续波”之间来回变形。更准确的说法是：不同实验分别显露出与两种经典图景相似的特征，但量子态本身不等同于任何一种经典对象。

\dialoguesupplement{教材在此排除的是朴素的经典确定路径模型：粒子预先选择一条不受另一条缝影响的路径，且两条路径的普通概率直接相加。仅凭本章实验，不应把结论扩大为排除所有带有轨迹概念的量子诠释。}

\notetopic{QM-C03-S05-T02}{从位置波函数到抽象状态及其动力学}

教材把粒子的状态写成位置波函数 $\psi(x,t)$，也写成抽象 ket $\ket{\psi(t)}$。二者不是两个状态，而满足
\begin{equation}
  \boxed{\psi(x,t)=\braket{x}{\psi(t)}}.
  \label{eq:QM-C03-S05-position-component}
\end{equation}
$\ket{\psi(t)}$ 是不依赖所选基的抽象态矢量；$\psi(x,t)$ 是它在连续位置基 $\{\ket{x}\}$ 下的坐标。利用完备关系
\begin{equation}
  \int_{-\infty}^{\infty}\ket{x}\bra{x}\dd x=I,
  \label{eq:QM-C03-S05-position-completeness}
\end{equation}
可以反向重建
\begin{align}
  \ket{\psi(t)}
  &=I\ket{\psi(t)} \notag\\
  &=\int_{-\infty}^{\infty}
    \ket{x}\braket{x}{\psi(t)}\dd x \notag\\
  &=\boxed{
    \int_{-\infty}^{\infty}
    \psi(x,t)\ket{x}\dd x}.
  \label{eq:QM-C03-S05-state-reconstruction}
\end{align}
这正是有限维展开 $\ket{v}=\sum_i\ket{i}\braket{i}{v}$ 的连续版本。

经典力学与量子力学的动力学问题可作如下对照：
\begin{center}
\begin{tabularx}{0.94\textwidth}{@{}>{\raggedright\arraybackslash}p{2.6cm}XX@{}}
  \toprule
  & 经典力学 & 量子力学的初步图景 \\
  \midrule
  状态 & 相空间点 $(x(t),p(t))$ & 态矢量 $\ket{\psi(t)}$ \\
  位置描述 & 确定位置 $x(t)$ & 振幅 $\psi(x,t)=\braket{x}{\psi(t)}$ \\
  位置预测 & 理想情况下由轨迹确定 & $\abs{\psi(x,t)}^2$ 给出概率密度 \\
  动力学问题 & 求 $(x(0),p(0))\mapsto(x(t),p(t))$ & 求 $\ket{\psi(0)}\mapsto\ket{\psi(t)}$ \\
  \bottomrule
\end{tabularx}
\end{center}

只知道某一组测量结果的概率通常不足以确定状态。例如对离散正交基 $\{\ket{1},\ket{2}\}$，
\begin{equation}
  \ket{\phi_+}=\frac{\ket{1}+\ket{2}}{\sqrt2},
  \qquad
  \ket{\phi_-}=\frac{\ket{1}-\ket{2}}{\sqrt2}
  \label{eq:QM-C03-S05-relative-phase-states}
\end{equation}
在该基下都给出 $P(1)=P(2)=1/2$，但相对相位不同，换到干涉基后即可区分。这说明概率振幅包含比单组概率更多的信息。

第三章的双缝实验提供了量子公设的物理动机，却不能单独推出量子理论的全部结构。下一章将进一步规定 $\ket{\psi(t)}$ 包含哪些可观测信息、测量如何产生结果，以及状态怎样随时间演化。

\textbf{离散概率与连续概率密度。}
若测量基是离散的 $\{\ket{n}\}$，则
\begin{equation}
  \boxed{P(n)=\abs{\braket{n}{\psi}}^2},
  \qquad
  \sum_nP(n)=1,
  \label{eq:QM-C03-S05-discrete-born-rule}
\end{equation}
这里的 $P(n)$ 是 probability mass（离散概率），不需要积分。若测量位置这一连续变量，则 $\abs{\psi(x)}^2$ 是 probability density（概率密度），有限区间概率必须按式~\eqref{eq:QM-C03-S05-position-born-rule} 积分。不能把连续空间中某一点的密度值直接当作该点的概率。

\textbf{一分钟回顾。}

量子对象在单次探测中以完整的局域事件出现，但其重复实验分布必须由概率振幅的叠加与 Born 规则计算；这就是本章所说波粒二象性的操作性内容。抽象态 $\ket{\psi(t)}$ 与位置波函数由 $\psi(x,t)=\braket{x}{\psi(t)}$ 联系，后者只是前者在位置基中的坐标。量子动力学研究的是整个状态 $\ket{\psi(t)}$ 如何演化，而不是预设一条经典确定轨迹。双缝实验为这些结构提供动机，但完整公设仍需由下一章正式给出。
